\documentclass[conference]{IEEEtran}
\usepackage{graphicx}
\usepackage{booktabs}
\usepackage{array}
\usepackage{url}

\title{Real Distributed AI Inference System with Cloud GPU Workers, Fault Tolerance, Monitoring, and Scalable Evaluation}

\author{
\IEEEauthorblockN{Boulos Aziz, [Team Member Names]}
\IEEEauthorblockA{CSE354 Distributed Computing\\
Ain Shams University, Faculty of Engineering\\
Spring 2025/2026}
}

\begin{document}
\maketitle

\begin{abstract}
This project implements and evaluates a real distributed AI inference system that runs large-language-model requests across multiple cloud GPU worker nodes. The system replaces a simulated distributed-computing prototype with a working deployment on six Thunder Compute NVIDIA RTX A6000 virtual machines running Ollama inference servers. A local controller performs scheduling, adaptive load balancing, request admission backpressure, health checking, retry logic, circuit-breaker-based fault tolerance, Retrieval-Augmented Generation (RAG), GPU telemetry collection, and formal evidence generation. The final evaluation tested 50, 100, 250, 500, and 1000 request workloads with increasing client concurrency. The six-worker system achieved 100\% success rate in all workloads, including the 1000-request/120-thread case, with zero failed requests, zero admission timeouts, all workers healthy, and GPU telemetry collected from all six nodes.
\end{abstract}

\begin{IEEEkeywords}
distributed computing, cloud GPU, AI inference, load balancing, fault tolerance, monitoring, RAG, Ollama, Thunder Compute, scalability
\end{IEEEkeywords}

\section{Introduction}
The CSE354 project specification requires a group project with a written report and developed code that demonstrates detailed analysis, design, testing, problem description, solution, limitations, sample output, references, and additional project information. The assessed learning outcomes include designing a distributed computing model, implementing it, configuring a working distributed environment, and communicating the work effectively.

The selected problem is distributed AI inference. A single inference server can become a bottleneck, fail under load, or lose availability during node failure. Therefore, a scalable inference service must distribute requests across multiple workers, detect failures, route around unhealthy nodes, measure performance, and produce evidence that the deployment is real rather than simulated.

\section{Problem Statement and Objectives}
The original project was an advanced distributed simulation. The target was to transform it into a real distributed AI inference system with cloud GPU workers, fault tolerance, monitoring, and scalable deployment. The main objectives were to replace simulated workers with real cloud GPU nodes, support concurrent users, detect and tolerate worker failures, balance load using live worker state, monitor GPU utilization and memory, and demonstrate a credible 1000-user evaluation.

\section{System Architecture}
The system follows a controller-worker architecture. Clients submit concurrent requests to a local controller. The controller builds retrieval context using a FAISS-based RAG pipeline, then dispatches each request to a remote Thunder GPU worker through an adaptive load balancer. Each GPU worker adapter communicates with an Ollama server over HTTPS using the OpenAI-compatible chat-completions API. A separate GPU metrics agent runs on each VM and exposes telemetry on port 9100.

% Convert SVG figures to PDF or PNG before compiling with IEEEtran.
% Recommended architecture figure: report/figures/architecture.svg

\section{Implementation}
\subsection{Real GPU Worker Deployment}
Workers are configured through environment variables, especially \texttt{LLM\_SERVER\_URLS}, \texttt{GPU\_METRICS\_URLS}, and \texttt{NUM\_WORKERS}. The final evaluation used six Thunder Compute workers: \texttt{ye96jt3q}, \texttt{a5oxns3p}, \texttt{uw01uuc2}, \texttt{xkj8xszu}, \texttt{hz8878v2}, and \texttt{ow2vbplc}.

\subsection{Fault Tolerance}
Each worker supports the states \texttt{HEALTHY}, \texttt{DEGRADED}, \texttt{UNHEALTHY}, and \texttt{RECOVERING}. Failures increment a worker failure counter; once the threshold is reached, the circuit breaker opens and the worker is excluded from dispatch. After cooldown and successful health checks, the worker recovers.

\subsection{Adaptive Load Balancing}
The load balancer scores workers using utilization, queue pressure, EWMA latency, oldest in-flight request age, failure rate, and state penalties. This allows the controller to favor faster workers while still using the full worker pool.

\subsection{Admission Backpressure}
At high concurrency, a worker pool may be healthy but temporarily saturated. The final system treats saturation as capacity pressure rather than node failure. Requests wait for capacity up to \texttt{SCHEDULER\_ADMISSION\_TIMEOUT}. In the 1000-request six-worker run, 126 requests waited for capacity, but zero timed out.

\subsection{Monitoring}
The GPU metrics agent queries \texttt{nvidia-smi} and exposes JSON metrics containing GPU utilization, memory usage, temperature, and power draw. The controller stores the samples in CSV and JSON evidence files.

\section{Experimental Setup}
The final deployment used six Thunder Compute prototyping VMs, each with one NVIDIA RTX A6000 GPU, four vCPUs, and 32 GB RAM. Each VM ran the Ollama template with TinyLlama. The formal evaluation matrix used 50/10, 100/20, 250/40, 500/80, and 1000/120 user/thread configurations.

\section{Results}
\begin{table}[h]
\caption{Six-Worker Evaluation Results}
\centering
\begin{tabular}{rrrrrr}
\toprule
Users & Success & Failed & Throughput & Avg. & P99 \\
 & (\%) & & (req/s) & (s) & (s) \\
\midrule
50 & 100.0 & 0 & 2.133 & 2.097 & 5.490 \\
100 & 100.0 & 0 & 3.778 & 2.895 & 9.844 \\
250 & 100.0 & 0 & 4.773 & 5.092 & 15.795 \\
500 & 100.0 & 0 & 5.686 & 9.791 & 58.728 \\
1000 & 100.0 & 0 & 6.905 & 13.931 & 115.463 \\
\bottomrule
\end{tabular}
\end{table}

\begin{table}[h]
\caption{Two-Worker Baseline vs. Six-Worker Final System}
\centering
\begin{tabular}{rrr}
\toprule
Users & Two Workers & Six Workers \\
\midrule
50 & 100.0\% & 100.0\% \\
100 & 100.0\% & 100.0\% \\
250 & 37.2\% & 100.0\% \\
500 & 9.8\% & 100.0\% \\
1000 & 5.0\% & 100.0\% \\
\bottomrule
\end{tabular}
\end{table}

The 1000-request run completed with 1000 total requests, zero failed requests, 100\% success rate, 6.905 requests per second, all workers healthy, 126 admission waits, and zero admission timeouts.

\section{Discussion}
The results show that the system scales horizontally. The two-worker system failed at high concurrency due to capacity exhaustion. The final system corrected this through both scale-out and admission backpressure. The adaptive load balancer favored faster workers while still using all six nodes, and the monitoring layer produced direct GPU evidence.

\section{Limitations and Future Work}
TinyLlama was selected for repeatable coursework evaluation; larger models would require more tuning. The 1000-request p99 latency is high because the system prioritizes successful completion through queueing rather than dropping requests. Future work includes Dockerization, authentication for public endpoints, persistent monitoring storage, autoscaling, and richer dashboard visualization.

\section{Conclusion}
The project successfully transformed a distributed-computing simulation into a real distributed AI inference system. The final implementation includes real Thunder Compute GPU workers, adaptive load balancing, fault tolerance, admission backpressure, RAG, GPU monitoring, and formal evidence. The six-worker deployment achieved 100\% success across the full evaluation matrix, including the 1000-request workload.

\begin{thebibliography}{00}
\bibitem{asu} Ain Shams University, Faculty of Engineering, ``CSE354: Distributed Computing Project Specification,'' Spring 2025/2026.
\bibitem{thunder} Thunder Compute, ``Use Instance Templates for AI,'' accessed May 10, 2026. [Online]. Available: \url{https://www.thundercompute.com/docs/guides/using-instance-templates}
\bibitem{ollama} Ollama, ``OpenAI Compatibility,'' accessed May 10, 2026. [Online]. Available: \url{https://docs.ollama.com/api/openai-compatibility}
\bibitem{faissdocs} Meta AI, ``Faiss Documentation,'' accessed May 10, 2026. [Online]. Available: \url{https://faiss.ai/index.html}
\bibitem{johnson} J. Johnson, M. Douze, and H. Jegou, ``Billion-Scale Similarity Search with GPUs,'' \emph{IEEE Transactions on Big Data}, vol. 7, no. 3, pp. 535--547, 2021, doi: 10.1109/TBDATA.2019.2921572.
\bibitem{metafaiss} Meta Engineering, ``Faiss: A library for efficient similarity search,'' 2017, accessed May 10, 2026. [Online]. Available: \url{https://engineering.fb.com/2017/03/29/data-infrastructure/faiss-a-library-for-efficient-similarity-search/}
\end{thebibliography}

\end{document}
